\documentclass{situationist-preprint}

\title{On the Aesthetic Value of Aesthetics}
\msid{AES-18-a}
\msfamily{AES-18}
\msdiscipline{Aesthetics}
\mstype{Principal manuscript}
\mskeywords{aesthetic theory; exemplification; theoretical virtues; elegance; performative consistency; Schlegel; Adorno}

\begin{document}
\maketitle

\begin{abstract}
Must aesthetic theory itself be beautiful? A tradition running from the Romantics to Adorno holds that a theory of art that does not share in the character of its object misunderstands it, and the thought has a simpler form: a theory that praises beauty in plain or ungainly prose seems to be at odds with itself. We argue that aesthetic theory is not required to exemplify the values it describes. A requirement on a discourse is grounded in its aim, and the aim of a theory is to describe and explain its subject truly. Exemplifying a value is neither necessary nor sufficient for describing or explaining it, as theories of colour, of humour and of grief show. Nor does a theory that commends beauty without having it commit any inconsistency, since it does not prescribe its own beauty. Aesthetic virtues of a theory, such as elegance, are permitted and may sometimes be evidence of its truth, but they are not required of it.
\end{abstract}

\section{The problem}

Friedrich Schlegel held, in one of his Critical Fragments, that poetry can be criticized only by way of poetry, and that a judgment of art which is not itself a work of art has no rights in the realm of art \citep{Schlegel1971}. \citet{Adorno1997} wrote his aesthetic theory in a paratactic form, without the hierarchy of premise and conclusion, on the ground that the form of a theory of art must answer to the form of its object. A plainer version of the thought is common: that a theory of beauty written without grace is somehow self-undermining, and that an aesthetics which is not itself aesthetically valuable betrays a failure to understand its subject.

We argue that aesthetic theory is not required to exemplify the values it describes. Following \citet{Goodman1968}, a thing \emph{exemplifies} a property when it both has the property and refers to it, as a tailor's swatch exemplifies its colour; we ask whether a theory of an aesthetic value must have, and so be able to exemplify, the value it is about.

\section{The argument}

\begin{principle}[Aim]\label{p:aim}
A requirement holds of a discourse as such only if meeting it is necessary for achieving the discourse's aim.
\end{principle}

\begin{principle}[Theory]\label{p:theory}
The aim of a theory of a subject is to describe and explain its subject truly, and so to make it understood.
\end{principle}

\begin{lemma}\label{l:indep}
Having an aesthetic value is neither necessary nor sufficient for describing or explaining it truly.
\end{lemma}

\begin{proof}
Not necessary: a theory of colour is not coloured, and it explains colour by describing the relations between surfaces, light and perceivers. A theory of humour explains why jokes are funny without being funny, and AES-13 implies that a successful explanation of a particular joke is not funny, since it resolves the incongruity on which the joke's effect depends. A theory of grief need not grieve. In each case what the theory must do is identify the features on which the value or property depends and the responses through which it is known, and nothing in that task requires that the theory have the value. Not sufficient: a beautiful text may describe its subject falsely, and many eloquent aesthetic theories, such as Hutcheson's identification of beauty with uniformity amidst variety, are elegant and false.
\end{proof}

\begin{theorem}\label{t:main}
Aesthetic theory is not required to exemplify the values it describes.
\end{theorem}

\begin{proof}
To exemplify a value, a theory must have it. By Principle~\ref{p:theory}, the aim of aesthetic theory is to describe and explain aesthetic values truly; by Lemma~\ref{l:indep}, having those values is not necessary for that aim. By Principle~\ref{p:aim}, having them is therefore not required of aesthetic theory as such, and neither is exemplifying them.
\end{proof}

Two clarifications. First, the theorem does not say that aesthetic virtues are irrelevant to theory. Elegance and simplicity are often taken as marks of good theories, and \citet{McAllister1996} argued that aesthetic criteria have played a real role in theory choice in the sciences; MTH-13 identifies mathematical beauty with the combination of surprise and inevitability, a property a theory may have. Where a theory's elegance is evidence of its truth, it is a reason to accept the theory, but its being so is a fact about the evidential connection, not a requirement that theories be beautiful. Second, clarity is required of a theory, since a theory that cannot be understood cannot make its subject understood; but clarity is a condition of communication, not an aesthetic value that the theory describes.

\section{Objections and replies}

\begin{objection}[Pragmatic inconsistency]\label{o:incons}
A theory that commends beauty while lacking it is at odds with itself, as a sermon on charity delivered with contempt would be.
\end{objection}

\begin{reply}
The sermon is inconsistent because it prescribes charity to all, its speaker included, and its delivery flouts the prescription. An aesthetic theory describes the conditions under which things are beautiful and why beauty is valuable; it does not prescribe that every text be beautiful, and so does not prescribe its own beauty. There is no prescription for its form to flout. The appearance of inconsistency comes from reading a theory of a value as a recommendation to have it.
\end{reply}

\begin{objection}[Acquaintance]\label{o:acq}
Aesthetic values are known only by acquaintance. A theory that does not give its reader the experience of the value cannot make it understood.
\end{objection}

\begin{reply}
A theory must presuppose acquaintance, not supply it. A theory of colour is understood by those who have seen colours, and a theory of musical expression by those who have heard music; it refers to the experiences and explains them, and the reader brings them. A theory that tried to supply the experience by exemplifying it would offer one instance of the value, which is not what a theory of the value is for.
\end{reply}

\begin{objection}[Romantic and critical theory]\label{o:adorno}
The object of aesthetics resists conceptual capture, and a theory that adopts the discursive form of other sciences falsifies it. Only a theory that shares the form of art can be true to art \citep{Adorno1997,Schlegel1971}.
\end{objection}

\begin{reply}
If art resisted every conceptual description, a theory in the form of art would not describe it either; it would be a further work. If art can be described, the description need not share its form. The objection is best taken as a claim about the limits of theory, that some of what matters in art is not captured by any description, and that claim is compatible with Theorem~\ref{t:main}, since it does not say that theory must be beautiful but that theory is not everything.
\end{reply}

\section{Conclusion}

Aesthetic theory is not required to be beautiful. Its aim is to describe and explain aesthetic values truly, and having those values is neither necessary nor sufficient for that aim. A theory that commends beauty without having it is not inconsistent, since it does not prescribe its own beauty, and it presupposes rather than supplies the reader's acquaintance with its subject. Elegance in a theory is permitted, and may be evidence, but it is not owed.

\begin{thebibliography}{99}\small
\bibitem[Adorno(1997)]{Adorno1997} Adorno, T.~W. (1997). \emph{Aesthetic Theory}. Trans. R. Hullot-Kentor. University of Minnesota Press. (Original work published 1970.)
\bibitem[Goodman(1968)]{Goodman1968} Goodman, N. (1968). \emph{Languages of Art: An Approach to a Theory of Symbols}. Bobbs-Merrill.
\bibitem[McAllister(1996)]{McAllister1996} McAllister, J.~W. (1996). \emph{Beauty and Revolution in Science}. Cornell University Press.
\bibitem[Schlegel(1971)]{Schlegel1971} Schlegel, F. (1971). \emph{Friedrich Schlegel's Lucinde and the Fragments}. Trans. P. Firchow. University of Minnesota Press. (Critical Fragments originally published 1797.)
\end{thebibliography}

\end{document}
